% !PRINTED 386
\weberpage{386}
%<<CONT
wenn $a$ \gesperrt{quadratischer Nichtrest} von $l$ ist, und nach der von
\gesperrt{Jacobi} eingeführten erweiterten Definition dieses Symbols ist
\weq{2}{\left(\frac{a}{l\,l'\,l''\dots}\right) = \left(\frac{a}{l}\right)\left(\frac{a}{l'}\right)\left(\frac{a}{l''}\right)\cdots}

Ist ferner $a$ ungerade, so ist
\weq{3}{\left(\frac{-1}{a}\right) = (-1)^{\frac{a-1}{2}}, \quad \left(\frac{-2}{a}\right) = (-1)^{\frac{a^2-1}{8}},}
und daraus, wenn $a, b, \dots$ relativ prim zu $m$ sind:
\weq{4}{\left(\frac{a}{m}\right)\left(\frac{b}{m}\right)\cdots = \left(\frac{ab\dots}{m}\right);}
ist $D$ eine Diskriminante und $a$ eine zu $2D$ teilerfremde Zahl,
die durch eine Form $\varphi$ der Diskriminante $D$ eigentlich darstellbar
 ist, ferner $l$ eine in $D$ aufgehende ungerade Primzahl,
so ist der Wert des Symbols $\left(\frac{a}{l}\right)$ nicht von der besonderen
Zahl $a$, sondern nur von der Formenklasse $A$, zu der $\varphi$ gehört,
abhängig. Dasselbe gilt von
\weq{5}{
\begin{gathered}
\left(\frac{-1}{a}\right) \quad D \equiv 0,-4 \pmod{16},\\
\left.\begin{aligned}
\left(\frac{2}{a}\right)& \quad D \equiv 0,8\\
\left(\frac{-2}{a}\right)& \quad D \equiv 0,-8
\end{aligned}\right\} \pmod{32}.
\end{gathered}
}

Hieraus lassen sich die Charaktere der Formenklassen
zusammensetzen.

Wir werden bei einer späteren Anwendung den Fall besonders
 zu berücksichtigen haben, daß $D$ durch 4 teilbar ist.
Wir setzen daher:
\weq{6}{D = -4m\footnotemark,}
\footnotetext{Nach der Bezeichnung von \gesperrt{Gauss: Formen erster Art der Determinante $-m$.}}
und erhalten für diesen Fall, wenn $l, l', \dots$ die verschiedenen
in $m$ aufgehenden ungeraden Primzahlen sind, die folgenden
Charaktere:
%CONT>>
